\NativeSlot{41}{36}{%
% p050_s01

\ExamQuestion{8.}{5}%
一个三位小数“四舍五入”后得到的近似数是7.20，这个三位小数最大是\mbox{（\hspace{1.8em}）}，最小是\mbox{（\hspace{1.8em}）}。\par
}
\NativeSlot{94.0}{129}{%
\ExamGroup{二、选择（将正确答案的序号填在括号里）。{\normalfont\kaishu\fontsize{10.8bp}{14bp}\selectfont （每题2分，共14分）}}

% p050_s02

\ExamQuestion{1.}{5}%
下图中不能表示0.06的是\mbox{（\hspace{1.8em}）}。\par
\ExamOptions{A. 图A}{B. 图B}{C. 图C}{D. 图D}
\BookFigure{figures/p050_s02-01.tex}
}
\NativeSlot{307.5}{36}{%
% p050_s03

\ExamQuestion{2.}{5}%
下面的数去掉0以后，大小不变的是\mbox{（\hspace{1.8em}）}。\par
\ExamOptions{A. 30.025}{B. 32.500}{C. 30025}{D. 300.25}
}
\NativeSlot{357.5}{36}{%
% p050_s04

\ExamQuestion{3.}{5}%
下面的小数中，与9最接近的是\mbox{（\hspace{1.8em}）}。\par
\ExamOptions{A. 8.9}{B. 9.01}{C. 9.001}{D. 8.89}
}
\NativeSlot{408.0}{92}{%
% p050_s05

\ExamQuestion{4.}{5}%
下图是某小学运动会3名选手的跳远情况，方框里数是\mbox{（\hspace{1.8em}）}。\par
\ExamOptions{A. 2}{B. 3}{C. 4}{D. 5}
\BookFigure{figures/p050_s05-01.tex}
}
\NativeSlot{564}{36}{%
% p050_s06

\ExamQuestion{5.}{5}%
\(0.\square8<0.6\)，方框里可以填的数有\mbox{（\hspace{1.8em}）}个。\par
\ExamOptions{A. 4}{B. 5}{C. 6}{D. 7}
}
\NativeSlot{617.5}{44}{%
% p050_s07

\ExamQuestion{6.}{5}%
用6、9、2、0这四个数字和小数点组成不同的三位小数，精确到个位是3的小数有\mbox{（\hspace{1.8em}）}个。\par
\ExamOptions{A. 6}{B. 5}{C. 4}{D. 3}
}
\NativeSlot{692.5}{49}{%
% p050_s08

\ExamQuestion{7.}{5}%
一个小数先扩大到它的100倍，再把小数点向左移动一位后是39.25，这个小数原来是\mbox{（\hspace{1.8em}）}。\par
\ExamOptions{A. 392.5}{B. 39.25}{C. 3.925}{D. 0.3925}
}
